%============================================================================
% LICENSING NOTICE
%============================================================================
% This WIP file, upon integration into guide-v.tex, becomes part of the
% TMS/DMS/CMS Usage Guide and is released under CC BY-NC 4.0.
% For details, see LICENSE in the repository root.
%============================================================================

%============================================================================
% FALCON BMS TMS/DMS/CMS HOTAS GUIDE
% WIP FILE TEMPLATE V1.0 — FINAL (Briefing v0.2.0.1)
%============================================================================

% IMPORTANTE: Este é um TEMPLATE padronizado para arquivos WIP (Work-In-Progress)
% que serão integrados ao guide.tex.

% Nomenclatura: Siga wip-naming-v1.4
% Padrão: section-C{N}-S{M}[-S{K}]-{titulo}-{status}-{data}.tex
% Exemplo: section-C5-S2-cms-actuation-hotas-tables-final-2026-01-10.tex

% Status: dev | review | final | approved | deprecated
% Locação: WIP/ (ativo) | ARCHIVE/ (aprovado/descartado)

%============================================================================
% PREAMBLE COMPLETO — TMS/DMS/CMS Usage Guide for Falcon BMS 4.38.1
% Gerado: 11 February 2026
% Status: Improvements as per GitHub Issues #34 and #33
%============================================================================

\documentclass[11pt, a4paper, twoside]{report}

% --------------------------------------------------------------------------
% BASIC ENCODING AND LANGUAGE
% --------------------------------------------------------------------------

\usepackage[utf8]{inputenc}
\usepackage[T1]{fontenc}
\usepackage[english]{babel}

% --------------------------------------------------------------------------
% FONTS AND MICROTYPOGRAPHY
% --------------------------------------------------------------------------

\usepackage{lmodern}
\usepackage{microtype}

% --------------------------------------------------------------------------
% PAGE GEOMETRY AND LAYOUT
% --------------------------------------------------------------------------

\usepackage{geometry}
\geometry{a4paper, left=2.0cm, right=2.0cm, top=2.5cm, bottom=2.5cm}
\usepackage{setspace}
\onehalfspacing

% --------------------------------------------------------------------------
% COLORS AND LINKS AND PDF SETTINGS
% --------------------------------------------------------------------------

\usepackage[table]{xcolor}
\definecolor{linkblue}{HTML}{004488}
\definecolor{linkred}{HTML}{882222}
\definecolor{headerblue}{HTML}{003366}
\definecolor{rowgray}{HTML}{F5F5F5}
\definecolor{subheadgray}{HTML}{E0E0E0}
\usepackage{soul}
\usepackage[pdfencoding=auto, psdextra, colorlinks=true, linkcolor=linkblue, citecolor=linkred, urlcolor=linkblue, breaklinks=true]{hyperref}
\hypersetup{
	pdftitle={TMS, DMS and CMS Usage Guide for Falcon BMS},
	pdfauthor={Carlos "Metal" Nader},
	pdfsubject={Flight Simulation - Falcon BMS HOTAS Reference},
	pdfkeywords={Falcon BMS, F-16, HOTAS, TMS, DMS, CMS, Flight Simulation},
	pdfcreator={pdfLaTeX (MiKTeX)},
	pdfproducer={MiKTeX},
	pdflang={en-US},
}
\usepackage{bookmark}

% --------------------------------------------------------------------------
% CAPTION SETUP BY TYPE
% --------------------------------------------------------------------------

\usepackage{caption}

% GLOBAL PATTERN
\captionsetup{font=small, labelfont=bf, justification=centering, singlelinecheck=true}
% FIGURES
\captionsetup[figure]{font=footnotesize, labelfont=bf, justification=centering, singlelinecheck=true}
% % TABLE/LONGTABLE/hotastable:
\captionsetup[table]{font=small, labelfont=bf, justification=centering, singlelinecheck=true}

% --------------------------------------------------------------------------
% HEADERS AND FOOTERS (IMPROVED for report + twoside)
% --------------------------------------------------------------------------

\usepackage{fancyhdr}
\setlength{\headheight}{25pt}                    
\pagestyle{fancy}
\fancyhf{}                                        
\fancyhead[LO,RE]{\small\textit{\leftmark}}     
\fancyhead[RO,LE]{\small\thepage}               
\fancyfoot{}                                      
\renewcommand{\headrulewidth}{0.4pt}
\renewcommand{\footrulewidth}{0pt}

% --------------------------------------------------------------------------
% CHAPTER FORMATTING AND SPACING (via titlesec)
% --------------------------------------------------------------------------

\usepackage{titlesec}

% --------------------------------------------------------------------------
% CHAPTER - Nível 0 (Destaque Máximo)
% --------------------------------------------------------------------------
% Shape: display → standalone em nova linha
% Font: \Large\bfseries --- label e título ambos em bold para destaque máximo
% Numeração: SIM
% Espaçamento: 15pt antes, 28pt depois
% --------------------------------------------------------------------------

\titleformat{\chapter}[display]
{\normalfont\Large\bfseries}           % Formato: LARGE + bold (label)
{\chaptertitlename~\thechapter}        % Rótulo: "Chapter 1" (bold também)
{20pt}                                 % Espaço vertical antes do corpo
{\Large\bfseries}                      % Corpo do título: LARGE + bold

\titlespacing{\chapter}
{0pt}       % Margem esquerda (sem indent)
{15pt}      % Espaço ANTES do título do capítulo
{28pt}      % Espaço DEPOIS do título do capítulo (AUMENTADO)
[0pt]       % Margem direita

% --------------------------------------------------------------------------
% SECTION - Nível 1
% --------------------------------------------------------------------------
% Shape: hang → label à esquerda, corpo indentado (compacto)
% Font: \large\bfseries (mesmo destaque que chapter para coerência)
% Numeração: SIM
% Espaçamento: 15pt antes, 8pt depois (REDUZIDO de 15pt/10pt anteriores)
% --------------------------------------------------------------------------

\titleformat{\section}[hang]
{\normalfont\large\bfseries}           % Formato: large + bold
{\thesection}                          % Rótulo: "1", "2", etc
{1em}                                  % Espaço entre rótulo e corpo
{}                                     % Corpo do título (herda do format)

\titlespacing{\section}
{0pt}       % Margem esquerda (sem indent)
{15pt}      % Espaço ANTES do título da seção (REDUZIDO)
{8pt}       % Espaço DEPOIS do título da seção (REDUZIDO)
[0pt]       % Margem direita

% --------------------------------------------------------------------------
% SUBSECTION - Nível 2
% --------------------------------------------------------------------------
% Shape: hang → label esquerda, corpo indentado
% Font: \normalsize\bfseries (redução visual, ainda robusto)
% Numeração: SIM
% Espaçamento: 12pt antes, 6pt depois (compacto)
% --------------------------------------------------------------------------

\titleformat{\subsection}[hang]
{\normalfont\normalsize\bfseries}      % Formato: tamanho normal + bold
{\thesubsection}                       % Rótulo: "1.1", "1.2", etc
{1em}                                  % Espaço entre rótulo e corpo
{}                                     % Corpo do título

\titlespacing{\subsection}
{0pt}       % Margem esquerda (sem indent)
{12pt}      % Espaço ANTES do título da subseção
{6pt}       % Espaço DEPOIS do título da subseção
[0pt]       % Margem direita

% --------------------------------------------------------------------------
% SUBSUBSECTION - Nível 3
% --------------------------------------------------------------------------
% Shape: hang → label esquerda, máximo recuo (bem compacto)
% Font: \normalsize\bfseries (mesmo tamanho que subsection, padrão)
% Numeração: SIM (configurable com secnumdepth)
% Espaçamento: 10pt antes, 4pt depois (bem comprimido)
% --------------------------------------------------------------------------

\titleformat{\subsubsection}[hang]
{\normalfont\normalsize\bfseries}      % Formato: tamanho normal + bold
{\thesubsubsection}                    % Rótulo: "1.1.1", "1.1.2", etc
{1em}                                  % Espaço entre rótulo e corpo
{}                                     % Corpo do título

\titlespacing{\subsubsection}
{0pt}       % Margem esquerda (sem indent)
{10pt}       % Espaço ANTES do título
{4pt}       % Espaço DEPOIS do título (bem comprimido)
[0pt]       % Margem direita

% --------------------------------------------------------------------------
% PARAGRAPH - Nível 4 (CRÍTICO: Ilusão Óptica Resolvida)
% --------------------------------------------------------------------------
% Shape: runin → inline, integrado ao texto (seu requisito)
% Font: \small\bfseries (SOLUÇÃO PARA ILUSÃO ÓPTICA DO NEGRITO)%       
% Numeração: NÃO (padrão para \paragraph)
% Espaçamento: runin (sem espaço vertical antes, integrado ao texto)
% --------------------------------------------------------------------------

\titleformat{\paragraph}[runin]
{\normalfont\small\bfseries}           % Formato: SMALL + bold (compensa ilusão)
{}                                     % Rótulo: vazio (sem numeração)
{0em}                                  % Espaço entre rótulo e corpo (zero, runin)
{}                                     % Corpo do título

\titlespacing{\paragraph}
{0pt}       % Margem esquerda (sem indent)
{8pt}       % Espaço ANTES do título do parágrafo (ADICIONADO - cria separação)
{1em}       % Espaço DEPOIS do título (ADICIONADO - espaço após ":")
[0pt]       % Margem direita

% --------------------------------------------------------------------------
% SUBPARAGRAPH - Nível 5 (Preparação Futura)
% --------------------------------------------------------------------------
% Shape: runin → inline, integrado ao texto
% Font: \small (MENOS destaque que \paragraph, sem bold)
% Numeração: NÃO (padrão para \subparagraph)
% Espaçamento: runin (integrado)
% Uso: Se usado, aparecerá com destaque menor que \paragraph
% --------------------------------------------------------------------------

\titleformat{\subparagraph}[runin]
{\normalfont\small}                    % Formato: small, SEM bold (menos destaque)
{}                                     % Rótulo: vazio (sem numeração)
{0em}                                  % Espaço entre rótulo e corpo (zero, runin)
{}                                     % Corpo do título

\titlespacing{\subparagraph}
{0pt}       % Margem esquerda (sem indent)
{8pt}       % Espaço ANTES do título (ADICIONADO)
{1em}       % Espaço DEPOIS do título (ADICIONADO)
[0pt]       % Margem direita

% --------------------------------------------------------------------------
% TABLES AND MACROS
% --------------------------------------------------------------------------

\usepackage{booktabs}
\usepackage{array}
\usepackage{longtable}
\usepackage{tabularx}

% Custom Columns
\newcolumntype{L}[1]{>{\raggedright\arraybackslash}p{#1}}
\newcolumntype{C}[1]{>{\centering\arraybackslash}p{#1}}
\newcolumntype{R}[1]{>{\raggedleft\arraybackslash}p{#1}}

% Macro for Visual Reference Links
\newcommand{\imglink}[1]{\hspace{2pt}\hyperref[#1]{\scriptsize\textbf{[Fig]}}}

% --------------------------------------------------------------------------
% HOTAS LOT FILE HANDLE 
% --------------------------------------------------------------------------
\makeatletter
\newcommand{\listofhotastables}{%
	\section*{List of HOTAS Tables}%
	\addcontentsline{toc}{section}{List of HOTAS Tables}%
	\@starttoc{hotas}%
}
\makeatother

% --------------------------------------------------------------------------
% HOTAS table environment
% Version: 2.1 (2026-02-10)
% --------------------------------------------------------------------------

\newenvironment{hotastable}[1]{%
	\footnotesize
	\setlength{\tabcolsep}{3pt}
	\renewcommand{\arraystretch}{1.35}
	\begin{longtable}{L{1.00cm} L{0.90cm} L{0.90cm} L{3.30cm} L{6.40cm} L{1.40cm} L{1.60cm}}
		\caption{#1}\label{table.\thetable}\\
		\noalign{\addtocontents{hotas}{\protect\contentsline{table}{\protect\numberline{\thetable}#1}{\thepage}{table.\thetable}}}%
		\rowcolor{headerblue}
		\multicolumn{1}{>{\centering\arraybackslash}p{1.00cm}}{\textbf{\color{white}Mode}} &
		\multicolumn{1}{>{\centering\arraybackslash}p{0.90cm}}{\textbf{\color{white}Dir.}} &
		\multicolumn{1}{>{\centering\arraybackslash}p{0.90cm}}{\textbf{\color{white}Act.}} &
		\multicolumn{1}{>{\centering\arraybackslash}p{3.30cm}}{\textbf{\color{white}Function}} &
		\multicolumn{1}{>{\centering\arraybackslash}p{6.40cm}}{\textbf{\color{white}Effect / Nuance}} &
		\multicolumn{1}{>{\centering\arraybackslash}p{1.40cm}}{\textbf{\color{white}Dash34}} &
		\multicolumn{1}{>{\centering\arraybackslash}p{1.60cm}}{\textbf{\color{white}Train.}} \\
		\endfirsthead
		\rowcolor{headerblue}
		\multicolumn{1}{>{\centering\arraybackslash}p{1.00cm}}{\textbf{\color{white}Mode}} &
		\multicolumn{1}{>{\centering\arraybackslash}p{0.90cm}}{\textbf{\color{white}Dir.}} &
		\multicolumn{1}{>{\centering\arraybackslash}p{0.90cm}}{\textbf{\color{white}Act.}} &
		\multicolumn{1}{>{\centering\arraybackslash}p{3.30cm}}{\textbf{\color{white}Function}} &
		\multicolumn{1}{>{\centering\arraybackslash}p{6.40cm}}{\textbf{\color{white}Effect / Nuance}} &
		\multicolumn{1}{>{\centering\arraybackslash}p{1.40cm}}{\textbf{\color{white}Dash34}} &
		\multicolumn{1}{>{\centering\arraybackslash}p{1.60cm}}{\textbf{\color{white}Train.}} \\
		\endhead
		\multicolumn{7}{r}{\small\emph{Continued on next page}}\\
		\endfoot
		\endlastfoot
	}{%
	\end{longtable}
}

% --------------------------------------------------------------------------
% SIMPLE REFERENCE MACROS FOR BMS DOCS
% --------------------------------------------------------------------------

\newcommand{\dashref}[1]{Dash-34~\S~#1}
\newcommand{\dashone}[1]{Dash-1~\S~#1}
\newcommand{\trnref}[1]{TRN~#1}
\newcommand{\trnman}{BMS Training Manual 4.38.1}
\newcommand{\bmsver}{Falcon BMS~4.38.1}
\newcommand{\dashrefs}[1]{\textit{TO 1F-16CMAM-34-1-1}, Dash-34, sections \texttt{#1}}

% --------------------------------------------------------------------------
% CROSS-REFERENCE MACROS (INTERNAL GUIDE)
% --------------------------------------------------------------------------

\newcommand{\secref}[1]{\hyperref[#1]{Section~\ref*{#1}}}
\newcommand{\chapref}[1]{\hyperref[#1]{Chapter~\ref*{#1}}}
\newcommand{\tabref}[1]{\hyperref[#1]{Table~\ref*{#1}}}
\newcommand{\figref}[1]{\hyperref[#1]{Figure~\ref*{#1}}}

% --------------------------------------------------------------------------
% VERSION CONTROL MACROS
% --------------------------------------------------------------------------

\newcommand{\docversion}{0.4.2.0}
\newcommand{\docbuild}{19 February 2026}
\newcommand{\docstartdate}{05 January 2026}
\newcommand{\docenddate}{19 February 2026}
\newcommand{\chapterscompletedof}{0/7}
\newcommand{\tablesfilledpct}{0}
\newcommand{\fulldocversion}{\docversion+\docbuild}

% --------------------------------------------------------------------------
% GRAPHICS
% --------------------------------------------------------------------------

\usepackage{graphicx}
\graphicspath{{fig/}}
\usepackage{float}
\usepackage{enumitem}

% --------------------------------------------------------------------------
% TITLE
% --------------------------------------------------------------------------

\title{TMS, DMS and CMS Usage Guide for \bmsver}
\author{Carlos ``Metal'' Nader}
\date{Version \fulldocversion{} | Progress: Chapters \chapterscompletedof{} | Tables \tablesfilledpct{} | January 2026}

% ============================================================================
% DOCUMENT BEGIN
% ============================================================================

\begin{document}
	
	\maketitle
	
	\pagenumbering{roman}
	
	% --------------------------------------------------------------------------
	% TOC DEPTH CONFIGURATION
	% --------------------------------------------------------------------------
	\setcounter{tocdepth}{3}       % Show up to \subsubsection in TOC
	\setcounter{secnumdepth}{3}    % Number up to \subsubsection
	
	\newpage
	\tableofcontents
	\newpage
	
	% --------------------------------------------------------------------------
	% LIST OF HOTAS TABLES
	% --------------------------------------------------------------------------
	\phantomsection
	\listofhotastables
	\newpage
	
	\pagenumbering{arabic}

%============================================================================
% WIP FILE METADATA (NOT RENDERED IN PDF)
%============================================================================

% File Name: chapter-C5-style-rev-review-2026-02-20.tex
% WIP Naming Convention: v1.4
% Target Chapter: C5 (CMS — Countermeasures Management Switch)
% Target Section: ALL (style revision pass — full chapter)
%
% WIP Status: review
% Created: 2026-02-19
% Last Modified: 2026-02-20
%
% Narrative Completion: ~40% (§5.1 completo, §5.2.1.1–§5.2.1.3 completos)
% Table Fill Status: N/A (style revision does not modify tables)
%
% Notes:
% - Contains approved revised prose for C5 only. Tables are not modified.
% - Each passage enters this file after explicit author approval.
% - Status remains 'review' until all C5 sections are complete; then → final.
% - Session 2026-02-19: §5.1 title revised, ¶1 rewritten, ¶2 and ¶3 eliminated.
%   §5.1.1 eliminated (content absorbed into §5.1).
% - Session 2026-02-20: §5.1 finalizado (figura e cross-references). §5.2 abertura,
%   lista de colunas, §5.2.1 (CMDS) e subseções §5.2.1.1–§5.2.1.3 completos.
% - Stopped at §5.2.2. Pending: §5.2.2 → §5.3.
% - TODO: RWR introduced for the first time at §5.2.1.2. Consider conceptual
%   introduction in C2 (HOTAS Fundamentals).
% - Cross-reference: Issue #44 (v0.4.2.0 style revision).

%============================================================================
% TEMPLATE STRUCTURE — Replace section/subsection headers with your content
%============================================================================

%============================================================================
% CONTENT BEGINS HERE (use \chapter{}, \section{}, \subsection{}, etc.)
%============================================================================
\chapter{CMS – Countermeasures Management Switch}
\label{chap:C5}

\section{Concept}
\label{sec:C5-S1}

The Countermeasures Management Switch (CMS) is a four-direction hat switch mounted on the control stick. It provides single-point control of the aircraft's defensive systems: the ALE-47 CMDS (Countermeasures Dispenser Set — the automatic chaff and flare dispensing system) and ECM (Electronic Countermeasures) systems, either external pods or internal avionics, depending on block configuration.

CMS actuation procedures are detailed in \secref{sec:C5-S2}. F-16 blocks and variant differences are covered in \secref{sec:C5-S3}.

For in-depth coverage of CMDS and ECM operation, see \dashrefs{2.7.1 and 2.7.2}, respectively. Diagrams of the flight stick and throttle controls are provided in \dashref{2.1.5}. \figref{fig:f16_hotas_cms_location} shows the CMS location on the F-16 control stick.

\begin{figure}[H]
\centering
\includegraphics[width=0.65\textwidth]{F-16_Side_Stick_Controller-1.jpg}
\caption{F-16 Throttle and Flight Stick. Image by Falconpedia (\url{falcon4.wikidot.com}), via Wikimedia Commons (\url{https://commons.wikimedia.org/wiki/File:F-16_Side_Stick_Controller.jpg}), licensed under the Creative Commons Attribution-Share Alike 3.0 Unported (CC BY-SA 3.0) license.}
\label{fig:f16_hotas_cms_location}
\end{figure}

\section{CMS Actuation}
\label{sec:C5-S2}

The tables in this section cover CMS usage by system: CMDS (Manual, Automatic, and Semi-Automatic modes) and ECM integration with both external pods (ALQ-131/ALQ-184) and internal avionics (IDIAS). 

CMS actuation is independent of the Master Mode currently selected.

Each table entry specifies:

\begin{itemize}
  \item \textbf{Mode}: The operational context (Master Mode).
  \item \textbf{Direction}: CMS hat direction (Up, Down, Left, Right).
  \item \textbf{Action}: Press type (Short, Long, Long Hold).
  \item \textbf{Function}: What the CMS command activates or controls.
  \item \textbf{Effect / Nuance}: Resulting system behavior, including tactics and constraints.
  \item \textbf{Dash34}: Reference section in the Dash-34 manual.
  \item \textbf{Training}: Recommended BMS training mission(s) (see \trnman{} for descriptions).
\end{itemize}

\subsection{CMDS}
\label{sec:C5-S2-S1}

The ALE-47 CMDS provides three operational modes: Manual (MAN), Automatic (AUTO), and Semi-Automatic (SEMI), each described below.

The CMS is the control interface for program execution and consent authority in all three modes.

\subsubsection{Manual Mode}
\label{sec:C5-S2-S1-S1}

In Manual (MAN) mode, the pilot executes countermeasure programs directly via CMS. Each CMS direction (\textbf{CMS Up} or \textbf{CMS Left}) executes a specific program.

\begin{hotastable}{CMS Actuation with CMDS Manual Mode}\label{tab:C5-S2-S1-S1-CMDS-Manual}
CMDS MAN & Up & Short & Execute Program 1--4 &
Runs the program selected via the CMDS panel PRGM knob once per press. No threat sensing; no automatic triggering. Overrides any AUTO dispensing, if running.
& 2.7.2.2 & 18, 28 \\
CMDS MAN & Left & Short & Execute Program 6 &
Flare-only program, activated regardless of PRGM knob position. Independent of threat sensing.
& 2.7.2.2 & 19 \\
\end{hotastable}

\subsubsection{Automatic Mode}
\label{sec:C5-S2-S1-S2}

In Automatic (AUTO) mode, the ALE-47 CMDS dispenses the selected program (1–4) continuously in response to threats detected by the Radar Warning Receiver (RWR).

Pilot consent is given once by pressing \textbf{CMS Down}; dispensing continues until \textbf{CMS Right} is pressed or expendables are exhausted.

\begin{hotastable}{CMS Actuation with CMDS Automatic Mode}\label{tab:C5-S2-S1-S2-CMDS-Auto}
CMDS AUTO & Up & Short & Execute Program 1--4 &
Runs the program selected via the CMDS panel PRGM knob once, interrupting any ongoing AUTO dispensing. After execution, AUTO resumes if the threat persists.
& 2.7.2.2 & 18, 28 \\
CMDS AUTO & Left & Short & Execute Program 6 &
Manual flare-only program, overrides AUTO. After execution, AUTO resumes.
& 2.7.2.2 & 19 \\
CMDS AUTO & Down & Short & Give Consent; Enable AUTO Dispensing &
Grants consent for AUTO dispensing. RWR-detected threats trigger automatic execution of the program selected via the CMDS panel PRGM knob. Dispensing continues until the threat clears or CMS Right is pressed. Consent state persists if the pilot switches to MAN mode; returning to AUTO resumes auto-dispensing if a threat is detected.
& 2.7.2.1 & 18, 28 \\
CMDS AUTO & Right & Short & Cancel Consent; Disable AUTO &
Removes CMDS consent and places the ALE-47 in Standby. Automatic dispensing halts immediately. CMS Down must be re-issued to resume AUTO operation.
& 2.7.2.1 & 18, 28 \\
\end{hotastable}

\subsubsection{Semi-Automatic Mode}
\label{sec:C5-S2-S1-S3}

In Semi-Automatic (SEMI) mode, the ALE-47 CMDS dispenses one program per RWR-detected threat, with pilot consent required for each event.

When a threat is detected, the CMDS control unit displays \texttt{DISPENSE RDY} and the system sounds a \texttt{COUNTER} voice message.

\begin{hotastable}{CMS Actuation with CMDS Semi-Automatic Mode}\label{tab:C5-S2-S1-S3-CMDS-Semi}
CMDS SEMI & Up & Short & Execute Program 1--4 &
Runs the program selected via the CMDS panel PRGM knob once, independent of RWR threat state. After execution, CMDS resumes SEMI monitoring.
& 2.7.2.2 & 18, 28 \\
CMDS SEMI & Left & Short & Execute Program 6 &
Manual flare-only program, independent of SEMI threat detection. After execution, CMDS resumes SEMI monitoring.
& 2.7.2.2 & 19 \\
CMDS SEMI & Down & Short & Give Consent; Dispense One Program &
Executes one instance of the selected program in response to a \texttt{COUNTER} prompt. If the threat persists or a new threat appears, CMDS issues \texttt{COUNTER} again. Consent state is tracked; switching to AUTO while consent is active resumes auto-dispensing immediately on the next threat.
& 2.7.2.1 & 18, 28 \\
CMDS SEMI & Right & Short & Cancel Consent; Return to Standby &
Removes SEMI consent. CMDS halts monitoring and returns to Standby. \texttt{COUNTER} messages cease.
& 2.7.2.1 & 18, 28 \\
\end{hotastable}

%============================================================================
% END OF SECTION
%============================================================================

\end{document}